% Source-aligned English for printed p.192.
\clearpage
\SourcePage{192}

\LayerTarget{tohoku:corollary:4.3.2:1}
\textsc{Corollary 1.} \textit{Suppose that $A$ is annihilated by a
two-sided ideal $I$ of $O$ such that $O/I$ is a semisimple ring $K$. Then the
two spectral sequences of the bicomplex $\mathrm{Hom}_O(L,\Gamma \mathrm C(A))$ are
trivial, and their initial terms identify canonically with the homology
$\mathrm{Ext}_{\mathbf O}(X;\mathbf M,A)$ of this bicomplex. More precisely,
there are canonical isomorphisms}

\LayerTarget{tohoku:eq:4.3.4}
\[
\tag{4.3.4}
\mathrm{Ext}_{\mathbf O}^n(X;\mathbf M,A)=
\bigoplus_{p+q=n}
\mathrm{Hom}_K(\mathrm{Tor}_p^O(K,M),H^q(X,A)).
\]
\LayerTarget{tohoku:proof:corollary-4.3.2:1}
Indeed, $\mathrm C(A)$, and hence $\Gamma \mathrm C(A)$, is also annihilated by $I$, whence

\LayerTarget{tohoku:eq:4.3.5}
\[
\tag{4.3.5}
\mathrm{Hom}_O(L,\Gamma \mathrm C(A))=
\mathrm{Hom}_K(K\otimes_O L,\Gamma \mathrm C(A))
\]
\LayerTarget{tohoku:proof:kunneth-collapse-4.3.4}
Since $K$ is semisimple, we may apply the simplest form of the Künneth theorem,
which follows here from the fact that $\mathrm{Hom}_K$ is an exact bifunctor.
This gives the triviality of the spectral sequences, together with the formula
$H(\mathrm{Hom}_K(K\otimes_O L,\Gamma \mathrm C(A)))=
\mathrm{Hom}_K(H(K\otimes_O L),H(\Gamma \mathrm C(A))),$
which is precisely (4.3.4).

\LayerTarget{tohoku:construction:relative-k-hyperhomology-spectral-sequences}
Continue to suppose that $A$ is annihilated by a two-sided ideal $I$, but make
no assumption on $K=O/I$. Equation (4.3.5) leads us to consider the
hyperhomology spectral sequences of the bifunctor $\mathrm{Hom}_K$ with
respect to the complexes $K\otimes_O L$ and $\Gamma \mathrm C(A)$. In the first, the
initial term

\LayerTarget{tohoku:formula:first-relative-k-initial-term}
$H^p(\mathrm{Ext}_K^q(K\otimes_O L,\Gamma \mathrm C(A)))$ vanishes for $q>0$, since
$K\otimes_O L$ is $K$-free. Thus the abutment identifies canonically with
$H(\mathrm{Hom}_K(K\otimes_O L,\Gamma \mathrm C(A)))$, which is precisely what must be
computed. In the second spectral sequence, the initial term is
\LayerTarget{tohoku:formula:second-relative-k-initial-term}
$\displaystyle \bigoplus_{q'+q''=q}
\mathrm{Ext}_K^p(H^{q'}(K\otimes_O L),H^{q''}\Gamma \mathrm C(A)).$
Hence:

\LayerTarget{tohoku:corollary:4.3.2:2}
\textsc{Corollary 2.} \textit{Suppose that the $\mathbf O$-module
$A$ is annihilated by a two-sided ideal $I$ of $O$, and put $K=O/I$. Then
$(\mathrm{Ext}_{\mathbf O}^n(X;\mathbf M,A))$ is also the abutment of a third
spectral functor, whose initial term is}

\LayerTarget{tohoku:eq:4.3.6}
\[
\tag{4.3.6}
\mathrm{III}_2^{pq}(A)=
\bigoplus_{q'+q''=q}
\mathrm{Ext}_K^p(\mathrm{Tor}_{q'}^O(K,M),H^{q''}(X,A)).
\]
\LayerTarget{tohoku:corollary:left-hereditary-k}
\textit{In particular, if $K$ is a \textup{``}left hereditary ring\textup{''} \textup{[6, 1.5]}, for
example a principal ideal domain, there is a canonical exact sequence}

\LayerTarget{tohoku:sequence:left-hereditary-k}
\[
\begin{split}
0\to{}&
\bigoplus_{p+q=n-1}
\mathrm{Ext}_K^1(\mathrm{Tor}_p^O(K,M),H^q(X,A))\\
\to{}&\mathrm{Ext}_{\mathbf O}^n(X;\mathbf M,A)
\to
\bigoplus_{p+q=n}
\mathrm{Hom}_K(\mathrm{Tor}_p^O(K,M),H^q(X,A))
\to0.
\end{split}
\]

\LayerTarget{tohoku:remark:4.3.2:1}
\textsc{Remarks.} 1. Suppose that the $L_i$ are isomorphic to
$O^{n_i}$. Then, by Proposition 4.3.2, the first hyperhomology spectral
sequence of the functor $\mathrm{Hom}_O(M,-)$ for the complex $\Gamma \mathrm C(A)$
is that of Theorem 4.2.1, whose initial term therefore has the remarkable
interpretation (4.2.5). Notice carefully, however, that this fact was deduced
from Corollary 3 to Proposition 4.2.4 and depends essentially on (4.2.11),
that is, \textit{on the fact that the canonical-resolution functor $\mathrm C$ commutes
with} $\mathbf{Hom}_{\mathbf O}(\mathbf L,-)$ when $L$ is an $O$-module isomorphic to $O^n$
(which, moreover, places us under the hypotheses of the corollary to
Proposition 2.5.2). This is

\TohokuContinuation{193}
